% Revised under the editorial brief; template formatting is unchanged.
\section{\src{Introduction and related literature}}\label{sec:introduction}
\subsection{\src{Motivation and financial relevance}}\label{sec:motivation}

\src{Many official lenders provide conditional financing: they disburse a loan in tranches and release each tranche after the first only when a review confirms that the borrower is meeting the policy conditions agreed at approval. The World Bank’s development policy financing and the European Union’s macro-financial assistance, for example, also attach policy conditions to their loans. We study interruption risk using programmes of the International Monetary Fund (IMF), which form a natural sample: the IMF’s Monitoring of Fund Arrangements database (MONA) records the review schedule of IMF-supported programmes approved since 2002 and whether each review was completed (IMF, 2026b), which gives a consistent record at the programme level.}

\src{The IMF supports countries facing balance-of-payments difficulties through lending programmes that combine financing with policy commitments and periodic reviews. We ask whether information available at approval by the IMF Executive Board (the Board) can predict whether a programme is later interrupted (see Section~\ref{sec:prediction-task} for the outcome definition).}

\src{Interruption matters financially because implementation can affect expected financing and confidence in the government’s policy commitments. Mody and Saravia (2006, pp. 843–845, 857) find that IMF programmes do not uniformly improve access to international bond markets; favourable effects depend on credible reform commitments and economic circumstances. Reinsberg et al. (2022, pp. 1023, 1034–1035) associate permanent interruptions with lower institutional investor ratings and a short-lived increase in sovereign bond spreads.}

\subsection{\src{Political incentives and implementation capacity}}\label{sec:political-incentives}

\src{The implementation literature emphasises the borrowing government’s willingness and ability to fulfil its commitments. Ivanova et al. (2003, pp. 4, 34) identify political cohesion, stability and bureaucratic quality as important influences on implementation. Joyce (2006, pp. 342–345, 351) similarly links implementation to government ideological cohesion, regime duration and political and economic openness. Dreher (2003, pp. 101, 111–114) shows that electoral incentives interact with political institutions and also identifies associations with initial economic vulnerability.}

\src{Our 13 inputs do not directly measure elections, political cohesion or bureaucratic quality, so these mechanisms are only partly observed.}

\subsection{\src{Conditionality and programme design}}\label{sec:conditionality}

\src{Evidence on conditionality is mixed. Ivanova et al. (2003, pp. 4, 34) find limited explanatory power for conditionality once domestic political factors are considered. Dreher’s (2009, pp. 248–250) survey discusses competing arguments about conditions as commitment devices and constraints on implementation. More recently, Reinsberg et al. (2022, pp. 1023, 1026, 1035) find that greater structural conditionality at programme onset is associated with interruption, especially conditions involving public-sector reform, privatisation and price liberalisation.}

\src{The outcome definition also affects interpretation, because duration and review counts may predict interruption partly through exposure (Section~\ref{sec:variables}). Counting conditions also misses their political difficulty and content.}

\subsection{\src{Macroeconomic vulnerability and prior programme use}}\label{sec:prior-use}

\src{Weak growth, inflation and external financing pressures can make agreed adjustment harder to implement. Iseringhausen et al. (2019, pp. 8, 16, 19, 23–25) find that macroeconomic conditions, institutions and previous IMF experience help explain repeated programme use. Their outcome, however, is a successor arrangement within three years. A country can return to the IMF after successfully completing a programme because it experiences a new shock.}

\subsection{\src{Machine learning in related work and our contribution}}\label{sec:contribution}

\src{Machine learning already contributes to IMF forecasting. Agbloyor et al. (2023, pp. 1–5) predict IMF bailouts using 6,550 observations and 138 features, reporting predictive information from economic, agricultural, energy, health and social indicators. Batsuuri et al. (2024, pp. 12, 14–18, 22–25) find that machine-learning methods outperform traditional approaches when forecasting new IMF arrangements, with results sensitive to preprocessing and prediction horizon. Both studies predict participation; we predict what happens after a programme is approved.}

\src{Gu et al. (2020, pp. 2223, 2225–2227) demonstrate the value of non-linearities and interactions in equity-return prediction. Their large stock-return panel differs from our 173-programme sample, so gains from flexibility may not carry over; logistic regression provides a benchmark. Breiman (2001a, pp. 5–8, 10–11) provides the basis for random forests, which combine predictions from randomised trees. We compare the random forest with logistic regression, an interpretable benchmark for a binary outcome, and a neural network.}

\src{Our contribution is a focused predictive test, in the sense of Shmueli (2010), of the research question stated in Section~\ref{sec:prediction-task}. The literature covers implementation mechanisms, financial consequences, predictors of participation or repeated use, and the determinants of interruption, but to our knowledge no study evaluates out-of-sample prediction of interruption with approval-day information.}
